% Generated draft: requires mathematical and source review.
\documentclass[11pt]{article}
\usepackage[margin=25mm]{geometry}
\usepackage{fontspec}
\setmainfont{Latin Modern Roman}
\usepackage{amsmath,amssymb,mathtools,unicode-math}
\setmathfont{Latin Modern Math}
\usepackage{enumitem,needspace,xurl,hyperref}
\setlength{\parindent}{0pt}
\setlength{\parskip}{5pt}
\setlength{\emergencystretch}{3em}
\newcommand{\N}{\mathbb{N}}
\newcommand{\Z}{\mathbb{Z}}
\newcommand{\Q}{\mathbb{Q}}
\newcommand{\R}{\mathbb{R}}
\newcommand{\C}{\mathbb{C}}
\newcommand{\F}{\mathbb{F}}
\newcommand{\A}{\mathbb{A}}
\newcommand{\PP}{\mathbb{P}}
\newcommand{\E}{\mathbb{E}}
\newcommand{\eps}{\varepsilon}
\newcommand{\dd}{\,\mathrm{d}}
\begin{document}
\section*{UnsolvedMath ID: 30005506}
\textbf{Source problem code:} catalogue entry
\subsection{Definitions and mathematical statement}
For a finite group \(K\), an \(\ell\)-block \(c\) of \(K\), and a defect group \(P\) of \(c\), let \(\operatorname{Irr}(c)\) denote the ordinary irreducible complex characters belonging to \(c\). For a positive integer \(m\), write \(m_\ell\) for its largest power-of-\(\ell\) divisor. Define the height-zero characters by
\[
\operatorname{Irr}_0(c)=\left\{\chi\in\operatorname{Irr}(c):\chi(1)_\ell=\frac{|K|_\ell}{|P|}\right\}.
\]
The general Alperin--McKay conjecture recalled in the source asserts that, for every finite group \(K\), every prime \(\ell\mid |K|\), and every \(\ell\)-block \(c\) with defect group \(P\),
\[
|\operatorname{Irr}_0(c)|=|\operatorname{Irr}_0(c^{\mathrm{Br}})|,
\]
where \(c^{\mathrm{Br}}\) is the Brauer correspondent of \(c\) in \(N_K(P)\).

The more specific construction requested by this record has the following supplied setup. Let \(\mathbf G\) be a simple algebraic group of simply connected type \(B\) or \(C\) over an algebraically closed field of characteristic \(p>0\). Let \(F:\mathbf G\to\mathbf G\) be a Frobenius endomorphism defining an \(\mathbb F_q\)-structure, put \(H=\mathbf G^F\), and let \(\ell\geq5\) be a prime distinct from \(p\). For every isolated \(\ell\)-block \(b\) of \(H\), the excerpt asserts the existence of a labeling \((\mathbf L,\lambda)\), where \(\mathbf L\) is a \(d\)-split Levi subgroup of \((\mathbf G,F)\) and \(\lambda\in\operatorname{Irr}(\mathbf L^F)\) is \(d\)-cuspidal. The excerpt does not specify \(d\) or define these specialized conditions; their intended meanings must be recovered from the cited report before this becomes a complete standalone formulation.

Let \(D\) be a defect group of \(b\) compatible with this labeling. The excerpt states that
\[
Z(\mathbf L)^F_\ell\mathrel{\mathrm{char}}D,
\]
where \(Z(\mathbf L)^F_\ell\) is the Sylow \(\ell\)-subgroup of the finite abelian group \(Z(\mathbf L)^F\), and \(\mathrel{\mathrm{char}}\) means characteristic subgroup. Under the intended finite rational-normalizer interpretation, set
\[
N=N_H(\mathbf L)=\{h\in H:h\mathbf Lh^{-1}=\mathbf L\},
\]
and let \(c_\lambda\) be the \(\ell\)-block of \(\mathbf L^F\) containing \(\lambda\). The excerpt asserts that there is a unique \(\ell\)-block \(B\) of \(N\) covering \(c_\lambda\). Here covering means that the restriction of some character in \(\operatorname{Irr}(B)\) has an irreducible constituent in \(\operatorname{Irr}(c_\lambda)\). This \(B\) is a block of a Levi normalizer; it must not be identified without further justification with the Brauer correspondent in \(N_H(D)\).

Put
\[
A=\operatorname{Aut}(H)_{b,D}
 =\{\alpha\in\operatorname{Aut}(H):\alpha(b)=b,\ \alpha(D)=D\}.
\]
The requested construction is an \(A\)-equivariant bijection
\[
\Phi_b:\operatorname{Irr}_0(b)\xrightarrow{\sim}\operatorname{Irr}_0(B),
\qquad
\Phi_b(\chi^\alpha)=\Phi_b(\chi)^\alpha
\quad(\chi\in\operatorname{Irr}_0(b),\ \alpha\in A).
\]
For an invariant finite group, the character action is \(\chi^\alpha(x)=\chi(\alpha^{-1}(x))\). On the local side this formula requires an action of \(A\) on \(N\) preserving \(B\), or a specified replacement action. The supplied excerpt does not explain that action or the necessary invariance, so the displayed target is conditional on resolving this point. The excerpt reports that this construction suffices for the inductive Alperin--McKay condition in its setting; it does not supply the full inductive condition or the hypotheses of that sufficiency reduction.

The cleaned statement names finite simple groups of Lie type \(B\) or \(C\), whereas the detailed excerpt works with \(H=\mathbf G^F\) for a simply connected algebraic group. These finite groups need not themselves be simple. The intended passage to simple central quotients, including any exceptional parameters, is unspecified. Also, the setup allows every \(p>0\), while the supplied dated literature assessment claims the stated cases are proved only with odd defining characteristic. Neither a restriction to odd \(p\) nor a resolution in characteristic two is silently imposed here.
\subsection{Short English statement}
For each isolated block in the supplied type \(B\) or \(C\) setting, construct a bijection between its height-zero characters and those of the specified Levi-normalizer block, compatible with the stated automorphism group. The source reports that this yields the relevant inductive Alperin--McKay condition. The local automorphism action, specialized block-labeling definitions, and passage from simply connected groups to finite simple groups require clarification. The supplied partially solved status is dated curation, not an independently verified conclusion.
\subsection{Sources}
\begin{itemize}
\item Supplied local UnsolvedMath registry record, permanent ID \(30005506\), exact source problem code \texttt{catalogue entry}, title: \emph{Inductive Alperin--McKay Condition for Isolated Blocks}. The registry marks the statement \texttt{corrected\_verified} and reports that it was reconstructed and checked against the official report; this is supplied provenance, not verification performed for this draft.
\item \emph{Representations of Finite Groups} (2023), Oberwolfach Reports, pp.~1031--1084, DOI \texttt{10.4171/owr/2023/19}: \url{https://doi.org/10.4171/owr/2023/19}. The supplied literature label is ``Ruhstorfer--Brough, inductive Alperin--McKay condition.'' The excerpt describes joint work in progress with Julian Brough; the record does not separately identify a proposer.
\item Supplied literature reference: ``Ruhstorfer, Alperin--McKay for the prime two,'' \url{https://annals.math.princeton.edu/2025/201-2/p02}.
\item Supplied literature reference: ``Ruhstorfer, Quasi-isolated blocks,'' \url{https://arxiv.org/abs/2103.06394}.
\item The excerpt cites [Sp\"a13] for the reduction to the inductive Alperin--McKay condition and [Ruh22] for the reduction to isolated blocks in types \(B\) and \(C\). Complete bibliographic entries and the precise scope of these reductions are not supplied.
\item Literature assessment date: \texttt{2026-08-22}. This is supplied dated curation, not independently verified. It labels the problem partially solved, states that the indicated type \(B\) and \(C\) cases with odd defining characteristic and \(\ell\geq5\) are proved, and states that the all-groups, all-primes conjecture remains open. The record itself requests independent verification of time-sensitive status and cited claims.
\end{itemize}
\end{document}
